\section{Local extrema and critical points}

A local maximum or local minimum is a place where the function is higher or lower than nearby values. Derivatives help us find candidates for such points, but they do not replace interpretation.

The main idea is this: if a smooth graph changes from increasing to decreasing, it has a local maximum. If it changes from decreasing to increasing, it has a local minimum.

\subsection*{Local maxima and minima}

\begin{definitionbox}
A function \(f\) has a local maximum at \(a\) if \(f(a)\) is greater than or equal to nearby values of the function.

A function \(f\) has a local minimum at \(a\) if \(f(a)\) is less than or equal to nearby values of the function.
\end{definitionbox}

The word local means nearby. A local maximum does not have to be the largest value of the function on its whole domain. It only has to be largest in a neighborhood of the point.

\subsection*{Critical points}

Local extrema often occur at points where the derivative is zero or undefined.

\begin{definitionbox}
A critical point of \(f\) is a point \(a\) in the domain of \(f\) such that
\[
f'(a)=0
\]
or \(f'(a)\) does not exist.
\end{definitionbox}

Critical points are candidates for local extrema. They must be checked. Not every critical point is an extremum.

\begin{examplebox}
For
\[
f(x)=x^2,
\]
we have
\[
f'(x)=2x.
\]
The derivative is zero at \(x=0\). The function decreases before \(0\) and increases after \(0\). Therefore \(x=0\) is a local minimum.
\end{examplebox}

\begin{examplebox}
For
\[
f(x)=x^3,
\]
we have
\[
f'(x)=3x^2.
\]
The derivative is zero at \(x=0\). But the function is increasing both before and after \(0\). Therefore \(x=0\) is a critical point, but not a local extremum.
\end{examplebox}

These examples show why finding \(f'(x)=0\) is only the beginning of the analysis.

\subsection*{The first derivative test}

The first derivative test uses the sign of \(f'\) around a critical point.

\begin{theorembox}
Let \(a\) be a critical point of \(f\).
\begin{itemize}
\item If \(f'\) changes from positive to negative at \(a\), then \(f\) has a local maximum at \(a\).
\item If \(f'\) changes from negative to positive at \(a\), then \(f\) has a local minimum at \(a\).
\item If \(f'\) does not change sign at \(a\), then \(a\) is not a local extremum by this test.
\end{itemize}
\end{theorembox}

\begin{examplebox}
Find the local extrema of
\[
f(x)=x^3-3x.
\]
We computed earlier that
\[
f'(x)=3(x-1)(x+1).
\]
The critical points are
\[
x=-1
\qquad\text{and}\qquad
x=1.
\]
The sign of \(f'\) is positive on \((-\infty,-1)\), negative on \((-1,1)\), and positive on \((1,\infty)\).

At \(x=-1\), the derivative changes from positive to negative, so there is a local maximum. Its value is
\[
f(-1)=(-1)^3-3(-1)=2.
\]
At \(x=1\), the derivative changes from negative to positive, so there is a local minimum. Its value is
\[
f(1)=1^3-3=-2.
\]
\end{examplebox}

The derivative sign tells us the shape of the graph: increasing, then decreasing, then increasing.

\subsection*{Endpoints}

On a closed interval, endpoints must be considered separately. A derivative test describes behavior inside an interval, but an endpoint can still be the largest or smallest value on the interval.

For example, on the interval \([0,2]\), the function \(f(x)=x^2\) has its smallest value at \(x=0\) and largest value at \(x=2\). The endpoint \(x=2\) is not found by solving \(f'(x)=0\), but it matters for the interval.

\begin{summarybox}
When looking for extrema, ask:
\begin{itemize}
\item What is the domain or interval being studied?
\item Where is \(f'(x)=0\)?
\item Where does \(f'(x)\) fail to exist while \(f\) is defined?
\item How does the sign of \(f'\) change around each critical point?
\item Are endpoints part of the problem?
\item What are the actual function values at the candidate points?
\end{itemize}
\end{summarybox}

Extrema are not only derivative facts. They are statements about function values near a point or on an interval.